\documentclass[aps,preprint]{revtex4}%
\usepackage{amssymb}
\usepackage{amsfonts}
\usepackage{amsmath}
\usepackage{graphicx}%
\setcounter{MaxMatrixCols}{30}
%TCIDATA{OutputFilter=latex2.dll}
%TCIDATA{Version=5.50.0.2960}
%TCIDATA{CSTFile=revtex4.cst}
%TCIDATA{Created=Wednesday, October 10, 2018 08:49:06}
%TCIDATA{LastRevised=Friday, October 12, 2018 10:48:30}
%TCIDATA{<META NAME="GraphicsSave" CONTENT="32">}
%TCIDATA{<META NAME="SaveForMode" CONTENT="1">}
%TCIDATA{BibliographyScheme=Manual}
%TCIDATA{<META NAME="DocumentShell" CONTENT="Articles\SW\REVTeX 4">}
%BeginMSIPreambleData
\providecommand{\U}[1]{\protect\rule{.1in}{.1in}}
%EndMSIPreambleData
\newtheorem{theorem}{Theorem}
\newtheorem{acknowledgement}[theorem]{Acknowledgement}
\newtheorem{algorithm}[theorem]{Algorithm}
\newtheorem{axiom}[theorem]{Axiom}
\newtheorem{claim}[theorem]{Claim}
\newtheorem{conclusion}[theorem]{Conclusion}
\newtheorem{condition}[theorem]{Condition}
\newtheorem{conjecture}[theorem]{Conjecture}
\newtheorem{corollary}[theorem]{Corollary}
\newtheorem{criterion}[theorem]{Criterion}
\newtheorem{definition}[theorem]{Definition}
\newtheorem{example}[theorem]{Example}
\newtheorem{exercise}[theorem]{Exercise}
\newtheorem{lemma}[theorem]{Lemma}
\newtheorem{notation}[theorem]{Notation}
\newtheorem{problem}[theorem]{Problem}
\newtheorem{proposition}[theorem]{Proposition}
\newtheorem{remark}[theorem]{Remark}
\newtheorem{solution}[theorem]{Solution}
\newtheorem{summary}[theorem]{Summary}
\newenvironment{proof}[1][Proof]{\noindent\textbf{#1.} }{\ \rule{0.5em}{0.5em}}
\begin{document}
\preprint{HEP/123-qed}
\title[Short title for running header]{REV\TeX4}
\author{First Author}
\affiliation{My Institution}
\author{Second Author}
\affiliation{My Institution}
\author{Third Author}
\affiliation{Other Institution}
\keywords{one two three}
\pacs{PACS number}

\begin{abstract}
Shell document for REV\TeX{} 4.

\end{abstract}
\volumeyear{year}
\volumenumber{number}
\issuenumber{number}
\eid{identifier}
\date[Date text]{date}
\received[Received text]{date}

\revised[Revised text]{date}

\accepted[Accepted text]{date}

\published[Published text]{date}

\startpage{101}
\endpage{102}
\maketitle
\tableofcontents


\section{Modular Theory}

\subsubsection{Properties of States}

Let $f$ be a \emph{faithful state} then the $c^{\ast}$-algebra representation
$\Pi_{f}$ and the Hilbert space representation $\Phi_{f}$ of $\mathcal{A}$ are
faithful i.e.
\begin{equation}
\ker\left\{  \Pi_{f}\right\}  =\ker\left\{  \Phi_{f}\right\}  =\left\{
0\right\}
\end{equation}
and thus $c^{\ast}$-algebra and Hilbert space isomorphisms onto $\Pi
_{f}\left(  \mathcal{A}\right)  $ and $\Phi_{f}\left(  \mathcal{A}\right)  $
respectively$.$ The $c^{\ast}$-algebra representation $\Pi_{f}$ can be
developed to a $w^{\ast}$-algebra representation by considering $\Pi
_{f}\left(  \mathcal{A}\right)  ^{\prime\prime}$ $\subseteq\mathcal{B}\left(
\Phi_{f}\left(  \mathcal{A}\right)  \right)  .$

(Note that $\Pi_{f}\left(  \mathcal{A}\right)  ^{\prime\prime}$is the
enveloping $w^{\ast}$-algebra of the $c^{\ast}$-algebra $\Pi_{f}\left(
\mathcal{A}\right)  $, i.e. $\Pi_{f}\left(  \mathcal{A}\right)  ^{\prime
\prime}\supseteq\Pi_{f}\left(  \mathcal{A}\right)  $ and $\Pi_{f}\left(
\mathcal{A}\right)  ^{\prime\prime\prime\prime}=\Pi_{f}\left(  \mathcal{A}%
\right)  ^{\prime\prime}$ shows that $\Pi_{f}\left(  \mathcal{A}\right)
^{\prime\prime}$ is a $w^{\ast}$-algebra. As $\Pi_{f}\left(  \mathcal{A}%
\right)  ^{\prime\prime}$ is a $w^{\ast}$-algebra acting on a Hilbert space it
is a Von Neumann algebra).

The vector $\Phi_{f}\left(  I_{\mathcal{A}}\right)  $ is seperating and cyclic
for $\Pi_{f}\left(  \mathcal{A}\right)  $ and $\Phi_{f}\left(  \mathcal{A}%
\right)  ,$ i.e.
\begin{equation}
\left\{  \Pi_{f}\left(  A\right)  \Phi_{f}\left(  I_{\mathcal{A}}\right)
=0\,\text{for any }A\in\mathcal{A}\right\}  \Leftrightarrow\left\{  \Pi
_{f}\left(  A\right)  =0\right\}  ,
\end{equation}
(This is equivalent to $\ker\left\{  \Phi_{f}\right\}  =\left\{  0\right\}
$)$,$ and
\begin{equation}
\mathcal{H}_{f}\equiv\Phi_{f}\left(  \mathcal{A}\right)  =\overline{\left\{
\left.  \Pi_{f}\left(  A\right)  \Phi_{f}\left(  I_{\mathcal{A}}\right)
\right\vert A\in\mathcal{A}\right\}  }.
\end{equation}
The expectation values wrt the state $f$ are given by
\begin{equation}
f\left(  A\right)  =\left\langle \left.  \Phi_{f}\left(  I_{\mathcal{A}%
}\right)  \right\vert \Pi_{f}\left(  A\right)  \Phi_{f}\left(  I_{\mathcal{A}%
}\right)  \right\rangle _{f}.
\end{equation}
thus as $f$ is faithful
\begin{equation}
\left(  \left\langle \left.  \Phi_{f}\left(  A\right)  \right\vert \Phi
_{f}\left(  A\right)  \right\rangle _{f}.=f\left(  A^{\dagger}A\right)
\right)  =0\Leftrightarrow\left(  A^{\dagger}A=0\right)  \Leftrightarrow
\left(  A=0\right)
\end{equation}
and
\begin{equation}
\left(  \left\langle \left.  \Phi_{f}\left(  I_{\mathcal{A}}\right)
\right\vert \Pi_{f}\left(  A\right)  ^{\dagger}\Pi_{f}\left(  A\right)
\Phi_{f}\left(  I_{\mathcal{A}}\right)  \right\rangle _{f}.=f\left(
A^{\dagger}A\right)  =0\,\right)  \Leftrightarrow\left(  A=0\right)  .
\end{equation}
States such that (\U{e308}_f\left(  I_\mathcal{A}\right)   \text{is seperating
and cyclic)}$\Leftrightarrow$( $f$ is seperating and cyclic) are called
\emph{modular states. }If the state\emph{ }$f$ is modular then the vector
$\Phi_{f}\left(  I_{\mathcal{A}}\right)  $ is modular for $\Pi_{f}\left(
\mathcal{A}\right)  ^{\prime\prime}$ and vice-versa$.$ As $\Phi_{f}\left(
I_{\mathcal{A}}\right)  $ is modular it is a faithful vector state in
$\Phi_{f}\left(  \mathcal{A}\right)  $, (though it might not be a normal state
for $\Pi_{f}\left(  \mathcal{A}\right)  ^{\prime\prime}),$ however this
neither implies nor is implied by the faithfulness of $f$ on $\mathcal{A}$,
(every normalized vector in $\Phi_{f}\left(  \mathcal{A}\right)  $ produces a
normal state \emph{in} $\mathcal{B}\left(  \Phi_{f}\left(  \mathcal{A}\right)
\right)  ).$ Modular states are of particular importance both mathematically
(Type $III$ representations of $w^{\ast}$-algebras) and physically as
\emph{KMS states are modular}.

A\emph{ }state $f$ is \emph{normal} if there exists a density operator
$D_{f}\in\Pi_{f}\left(  \mathcal{A}\right)  ^{\prime\prime}$ such that
\begin{equation}
f\left(  A\right)  =\operatorname{Tr}\left\{  \Pi_{f}\left(  A\right)
D_{f}\right\}  =f_{\Pi_{f}}\left(  \Pi_{f}\left(  A\right)  \right)  ,
\end{equation}
remember $\left\vert \Phi_{f}\left(  I_{\mathcal{A}}\right)  \right\rangle
\left\langle \Phi_{f}\left(  I_{\mathcal{A}}\right)  \right\vert $ \emph{does
not have to belong to} $\Pi_{f}\left(  \mathcal{A}\right)  ^{\prime\prime}$
and in fact cannot belong to it unless $\Pi_{f}\left(  \mathcal{A}\right)
^{\prime\prime}$ is Type $I.$ The support, $\operatorname{Supp}\left\{
f\right\}  ,$ of a normal state $f$ $\in\mathcal{A}_{\ast}$ is the smallest
orthogonal projection, $P_{f}\in\mathcal{A}^{\prime\prime}$, that has the
property
\begin{equation}
f\left(  P_{f}\right)  =f_{\Pi_{f}}\left(  \Pi_{f}\left(  A\right)  \right)
=\operatorname{Tr}\left\{  \Pi_{f}\left(  P_{f}\right)  D_{f}\right\}  =1.
\end{equation}
$f$ is faithful $\operatorname*{Iff}$
\begin{equation}
\operatorname{Supp}\left\{  f\right\}  =I_{\mathcal{A}}.
\end{equation}
One can show that
\begin{equation}
\operatorname{Supp}\left\{  f_{\Pi_{f}}\right\}  =P_{f}=P_{\overline{\Pi
_{f}\left(  \mathcal{A}\right)  ^{\prime\prime\prime}\Phi_{f}\left(
I_{\mathcal{A}}\right)  }}%
\end{equation}
$\operatorname*{Iff}$ $\Phi_{f}\left(  I_{\mathcal{A}}\right)  $ is cyclic for
$\Pi_{f}\left(  \mathcal{A}\right)  ^{\prime\prime\prime},$ where
$P_{\overline{\Pi_{f}\left(  \mathcal{A}\right)  ^{\prime\prime\prime}\Phi
_{f}\left(  I_{\mathcal{A}}\right)  }}$ is the orthogonal projector on the
closure of $\left\{  \Pi_{f}\left(  \mathcal{A}\right)  ^{\prime\prime\prime
}\Phi_{f}\left(  I_{\mathcal{A}}\right)  \right\}  $ and $\Phi_{f}\left(
I_{\mathcal{A}}\right)  $ is seperating for $\Pi_{f}\left(  \mathcal{A}%
\right)  ^{\prime\prime}.$

$\lceil$The polar decomposition for any \emph{linear} operator, $A$, is
\begin{equation}
\text{ }A=U\left\vert A\right\vert =U\left(  A^{\dagger}A\right)  ^{\frac
{1}{2}}=\left(  AA^{\dagger}\right)  ^{\frac{1}{2}}U^{\dagger},
\end{equation}
where $U$ is an \emph{isometry}, thus if
\begin{equation}
A^{\dagger}A=0
\end{equation}
then
\begin{equation}
A=0.
\end{equation}
The polar decomposition for any \emph{conjugate linear} operator, $A$, gives
\begin{equation}
\text{ }A=U\left\vert A\right\vert =U\left(  A^{\dagger}A\right)  ^{\frac
{1}{2}}=\left(  AA^{\dagger}\right)  ^{\frac{1}{2}}U^{\dagger},
\end{equation}
where $U$ is an \emph{anti isometry}.$\rfloor$

\subsection{Tomita-Takesaki Theory}

Consider the conjugate linear map %

\begin{equation}
S_{0f}:\mathcal{H}_{f}\rightarrow\mathcal{H}_{f},
\end{equation}
where
\begin{equation}
\mathcal{H}_{f}=\Phi_{f}\left(  \mathcal{A}\right)  ,
\end{equation}
given by%
\begin{equation}
S_{0f}\left(  T\Phi_{f}\left(  I_{\mathcal{A}}\right)  \right)  \rightarrow
T^{\dagger}\Phi_{f}\left(  I_{\mathcal{A}}\right)  ,\,\forall T\in\Pi
_{f}\left(  \mathcal{A}\right)  ^{\prime\prime}.
\end{equation}
$S_{0f}$ is conjugate linear, with domain $\left\{  \left.  T\Phi_{f}\left(
I_{\mathcal{A}}\right)  \right\vert T\in\Pi_{f}\left(  \mathcal{A}\right)
^{\prime\prime}\right\}  $ and may be unbounded, but closeable, with closure
denoted by $S_{f}.$The adjoint $S_{f}^{\dagger},$ which is densley defined, is
defined by
\begin{equation}
\left\langle \left.  S_{f}^{\dagger}\Phi_{f}\left(  A\right)  \right\vert
\Phi_{f}\left(  B\right)  \right\rangle _{f}=\left\langle \left.  \Phi
_{f}\left(  A\right)  \right\vert S_{f}\Phi_{f}\left(  B\right)  \right\rangle
_{f}\,\forall B\in\mathcal{A}.
\end{equation}


One can also define the conjugate linear map
\begin{equation}
F_{0f}:\mathcal{H}_{f}\rightarrow\mathcal{H}_{f},
\end{equation}
that acts on the commutant, $\Pi_{f}\left(  \mathcal{A}\right)  ^{\prime
\prime\prime},$ of $\Pi_{f}\left(  \mathcal{A}\right)  ^{\prime\prime}$ by%
\begin{equation}
F_{0f}\left(  T\Phi_{f}\left(  I_{\mathcal{A}}\right)  \right)  \rightarrow
T^{\dagger}\Phi_{f}\left(  I_{\mathcal{A}}\right)  ,\,\,\forall T\in\Pi
_{f}\left(  \mathcal{A}\right)  ^{\prime\prime\prime},
\end{equation}
with closure denoted by $F_{f}.$ The adjoint $F_{f}^{\dagger}$ is defined in
similar way to $S_{f}^{\dagger}.$

$S$ induces the map
\begin{align*}
\widehat{S} &  :\Pi_{f}\left(  \mathcal{A}\right)  \rightarrow\Pi_{f}\left(
\mathcal{A}\right)  \\
\widehat{S}\left(  \Pi_{f}\left(  A\right)  \right)   &  \rightarrow\Pi
_{f}\left(  A\right)  ^{\dagger}=\Pi_{f}\left(  A^{\dagger}\right)  \\
\widehat{S}\left(  \Pi_{f}\left(  A\right)  \right)  \Phi_{f}\left(
I_{\mathcal{A}}\right)   &  =\Pi_{f}\left(  A\right)  ^{\dagger}\Phi
_{f}\left(  I_{\mathcal{A}}\right)
\end{align*}
that belongs to $\mathcal{L}^{CL}\left(  \Pi_{f}\left(  A\right)  \right)
\equiv\mathcal{L}^{CL}\left(  \mathcal{A}_{f}\right)  .$

One can observe that
\begin{equation}
\mathcal{L}^{CL}\left(  \mathcal{H}_{f}\right)  \cong\mathcal{L}\left(
\mathcal{H}_{f}\right)
\end{equation}
as $c^{\ast}$-algebras.

The adjoint map in $\mathcal{L}\left(  \mathcal{H}_{f}\right)  $ can be
extended to $\mathcal{L}^{CL}\left(  \mathcal{H}_{f}\right)  $
\begin{equation}
\left\langle \left.  A^{\dagger}\Phi_{f}\left(  B_{1}\right)  \right\vert
\Phi_{f}\left(  B_{2}\right)  \right\rangle _{f}=\left\langle \left.  \Phi
_{f}\left(  B_{1}\right)  \right\vert A\Phi_{f}\left(  B_{2}\right)
\right\rangle _{f}\quad\forall\Phi_{f}\left(  B_{2}\right)  \in\Phi_{f}\left(
\mathcal{A}\right)  ,\Phi_{f}\left(  B_{1}\right)  \in\Phi_{f}\left(
\mathcal{A}\right)  ,A\in\mathcal{L}^{CL}\left(  \mathcal{H}_{f}\right)  .
\end{equation}


$\lceil$The operator $A^{\dagger}$ is also conjugate linear. As the state $f$
is faithful the adjoint operations in $\mathcal{A}$, $\mathcal{B}^{CL}\left(
\mathcal{H}_{f}\right)  $ and $\mathcal{B}\left(  \mathcal{H}_{f}\right)  $
are the same. $\rfloor$

Equipped with the inner product
\begin{equation}
\left\langle \left.  B_{1}\right\vert B_{2}\right\rangle _{f}=f\left(
B_{1}^{\dagger}B_{2}\right)  =\left\langle \left.  \Phi_{f}\left(
B_{1}\right)  \right\vert \Phi_{f}\left(  B_{2}\right)  \right\rangle _{f}%
\end{equation}
$\mathcal{A}$ becomes a Hilbert algebra $\mathcal{A}_{f}$ and one has the
$c^{\ast}$-Hilbert algebra isomorphism
\begin{equation}
\mathcal{A}_{f}\cong\Pi_{f}\left(  \mathcal{A}\right)
\end{equation}
and the Hilbert space isomorphism
\begin{equation}
\mathcal{A}_{f}\cong\Phi_{f}\left(  \mathcal{A}\right)  .
\end{equation}
Thus one can consider $\widehat{S_{f}}$ as acting on $\mathcal{A}_{f}$ induced
by its action on $\mathcal{H}_{f}$ to give
\begin{equation}
\widehat{S_{f}}\left(  A\right)  =A^{\dagger}.
\end{equation}
and thus
\begin{equation}
\widehat{S_{f}}^{2}\left(  A\right)  =A\quad\forall A\in\mathcal{A}.
\end{equation}
$\lceil$%
\begin{align}
\left\langle \left.  S_{f}\left(  A\right)  \left(  B_{1}\right)  \right\vert
B_{2}\right\rangle _{f} &  =\left\langle \left.  B_{1}\right\vert A\left(
B_{2}\right)  \right\rangle _{f}\,\forall B_{2}\in\mathcal{H}_{f},B_{1}%
\in\mathcal{H}_{f}\\
f\left(  \left(  S_{f}\left(  A\right)  \left(  B_{1}\right)  \right)
^{\dagger}B_{2}\right)   &  =f\left(  B_{1}^{\dagger}A\left(  B_{2}\right)
\right)  \forall B_{2}\in\mathcal{H}_{f},B_{1}\in\mathcal{H}_{f}%
\end{align}%
\begin{align}
\left\Vert S_{f}\right\Vert _{\mathcal{A}_{f}}  & =\sup\limits_{B\in
\mathcal{H}_{f}}\left\{  \frac{\left\langle \left.  B\right\vert \left[
S_{f}^{\dagger}S_{f}\right]  \left(  B\right)  \right\rangle _{f}%
}{\left\langle \left.  B\right\vert B\right\rangle _{f}}\right\}
=\sup\limits_{B\in\mathcal{H}_{f}}\left\{  \frac{f\left(  B^{\dagger}\left[
S_{f}^{\dagger}S_{f}\right]  \left(  B\right)  \right)  }{f\left(  B^{\dagger
}B\right)  }\right\}  \\
& =\sup\limits_{B\in\mathcal{H}_{f}}\left\{  \frac{f\left(  \left(
S_{f}\left(  B\right)  \right)  ^{\dagger}S_{f}\left(  B\right)  \right)
}{f\left(  B^{\dagger}B\right)  }\right\}  =\sup\limits_{B\in\mathcal{H}_{f}%
}\left\{  \frac{f\left(  BB^{\dagger}\right)  }{f\left(  B^{\dagger}B\right)
}\right\}
\end{align}
$\rfloor$

and
\begin{equation}
S_{f}\left(  A^{\dagger}\right)  =A.
\end{equation}


As
\begin{equation}
S_{f}\in\mathcal{B}^{CL}\left(  \mathcal{H}_{f}\right)  \cong\mathcal{B}%
^{CL}\left(  \mathcal{A}_{f}\right)  \cong\mathcal{B}\left(  \mathcal{H}%
_{f}\right)  \cong\mathcal{B}\left(  \mathcal{A}_{f}\right)  ,
\end{equation}
one has
\begin{equation}
S_{f}\left(  A\right)  ^{\dagger}=A.
\end{equation}
$\lceil$The inner product is given by
\begin{equation}
\left\langle \left.  B_{1}\right\vert B_{2}\right\rangle _{f}=\overline
{\left\langle \left.  B_{2}\right\vert B_{1}\right\rangle _{f}}=f\left(
B_{1}^{\dagger}B_{2}\right)  =\overline{f\left(  B_{2}^{\dagger}B_{1}\right)
}%
\end{equation}


The adjoint (star) map is given by
\begin{equation}
\left\langle \left.  S_{f}\left(  B_{1}\right)  \right\vert B_{2}\right\rangle
_{f}=\left\langle \left.  B_{1}^{\dagger}\right\vert B_{2}\right\rangle
_{f}=f\left(  B_{1}B_{2}\right)  =f\left(  \left(  S_{f}\left(  B_{1}\right)
\right)  ^{\dagger}B_{2}\right)  =\left\langle \left.  B_{1}\right\vert
S_{f}^{\dagger}\left(  B_{2}\right)  \right\rangle _{f}=\overline{f\left(
B_{2}^{\dagger}B_{1}^{\dagger}\right)  }%
\end{equation}%
\begin{equation}
\left\langle \left.  S_{f}^{\dagger}\left(  B_{1}\right)  \right\vert
B_{2}\right\rangle _{f}=f\left(  \left(  S_{f}^{\dagger}\left(  B_{1}\right)
\right)  ^{\dagger}B_{2}\right)  =\left\langle \left.  B_{1}\right\vert
S_{f}\left(  B_{2}\right)  \right\rangle _{f}=\left\langle \left.
B_{1}\right\vert B_{2}^{\dagger}\right\rangle _{f}=f\left(  B_{1}^{\dagger
}B_{2}^{\dagger}\right)  =\overline{f\left(  B_{2}B_{1}\right)  }%
\end{equation}%
\begin{equation}
\left\langle \left.  B_{1}\right\vert S_{f}^{\dagger}\left(  B_{2}\right)
\right\rangle _{f}=\left\langle \left.  S_{f}\left(  B_{1}\right)  \right\vert
B_{2}\right\rangle _{f}=\left\langle \left.  B_{1}^{\dagger}\right\vert
B_{2}\right\rangle _{f}=f\left(  B_{1}^{\dagger}S_{f}^{\dagger}\left(
B_{2}\right)  \right)  =f\left(  B_{1}B_{2}\right)  =\overline{f\left(
B_{2}^{\dagger}B_{1}^{\dagger}\right)  }%
\end{equation}%
\begin{equation}
\left\langle \left.  B_{1}\right\vert S_{f}\left(  B_{2}\right)  \right\rangle
_{f}=f\left(  B_{1}^{\dagger}S_{f}\left(  B_{2}\right)  \right)  =\left\langle
\left.  S_{f}^{\dagger}\left(  B_{1}\right)  \right\vert B_{2}\right\rangle
_{f}=\left\langle \left.  B_{1}\right\vert B_{2}^{\dagger}\right\rangle
_{f}=f\left(  B_{1}^{\dagger}B_{2}^{\dagger}\right)  =\overline{f\left(
B_{2}B_{1}\right)  }%
\end{equation}


The postive operator $SS^{\dagger}$ is given by
\begin{align}
\left\langle \left.  \left(  S^{\dagger}\left(  B_{1}\right)  \right)
^{\dagger}\right\vert B_{2}\right\rangle  &  =\left\langle \left.
SS^{\dagger}\left(  B_{1}\right)  \right\vert B_{2}\right\rangle
_{f}=\left\langle \left.  B_{1}\right\vert SS^{\dagger}\left(  B_{2}\right)
\right\rangle _{f}=\left\langle \left.  B_{1}\right\vert SS^{\dagger}\left(
B_{2}\right)  \right\rangle _{f}=\left\langle \left.  S^{\dagger}\left(
B_{1}\right)  \right\vert S^{\dagger}\left(  B_{2}\right)  \right\rangle
_{f}\\
&  =f\left(  \left(  S^{\dagger}\left(  B_{1}\right)  \right)  ^{\dagger
}S^{\dagger}\left(  B_{2}\right)  \right)  =f\left(  B_{1}^{\dagger}\left[
SS^{\dagger}\right]  \left(  B_{2}\right)  \right)  =f\left(  S^{\dagger
}\left(  B_{1}\right)  B_{2}\right)  =\overline{\left\langle \left.
B_{2}\right\vert \left(  S^{\dagger}\left(  B_{1}\right)  \right)  ^{\dagger
}\right\rangle }_{f}%
\end{align}
and the postive operator $S_{f}^{\dagger}S_{f}$ by
\begin{align}
&  \left\langle \left.  S_{f}^{\dagger}\left(  B_{1}^{\dagger}\right)
\right\vert B_{2}\right\rangle _{f}=\left\langle \left.  S_{f}^{\dagger}%
S_{f}\left(  B_{1}\right)  \right\vert B_{2}\right\rangle _{f}=\left\langle
\left.  S_{f}\left(  B_{1}\right)  \right\vert S_{f}\left(  B_{2}\right)
\right\rangle _{f}\\
&  =\left\langle \left.  B_{1}\right\vert S_{f}^{\dagger}S_{f}\left(
B_{2}\right)  \right\rangle _{f}=\left\langle \left.  B_{1}^{\dagger
}\right\vert S_{f}\left(  B_{2}\right)  \right\rangle _{f}=f\left(  B_{1}%
B_{2}^{\dagger}\right)  =\left\langle \left.  B_{1}^{\dagger}\right\vert
B_{2}^{\dagger}\right\rangle _{f}.=f\left(  B_{1}^{\dagger}\left[
S_{f}^{\dagger}S_{f}\right]  \left(  B_{2}\right)  \right)  \ =\overline
{f\left(  B_{2}B_{1}^{\dagger}\right)  }%
\end{align}


The expectation values of the adjoint map are given by
\begin{align}
\left\langle \left.  S_{f}^{\dagger}\left(  B_{1}\right)  \right\vert
I\right\rangle _{f}  &  =\left\langle \left.  B_{1}\right\vert S_{f}\left(
I\right)  \right\rangle _{f}=\left\langle \left.  B_{1}\right\vert
I\right\rangle _{f}=f\left(  B_{1}^{\dagger}\right)  =\overline{f\left(
B_{1}\right)  }\\
\left\langle \left.  I\right\vert S_{f}^{\dagger}\left(  B_{1}\right)
\right\rangle _{f}  &  =\left\langle \left.  S_{f}^{\dagger}\left(  I\right)
\right\vert B_{1}\right\rangle _{f}=\left\langle \left.  I\right\vert
B_{1}\right\rangle _{f}=f\left(  B_{1}\right) \\
\left\langle \left.  S_{f}\left(  B_{1}\right)  \right\vert I\right\rangle
_{f}  &  =\left\langle \left.  B_{1}\right\vert S_{f}^{\dagger}\left(
I\right)  \right\rangle _{f}=\left\langle \left.  B_{1}\right\vert
I\right\rangle _{f}=f\left(  B_{1}^{\dagger}\right)  =\overline{f\left(
B_{1}\right)  }\\
\left\langle \left.  I\right\vert S_{f}\left(  B_{1}\right)  \right\rangle
_{f}  &  =\left\langle \left.  S_{f}^{\dagger}\left(  I\right)  \right\vert
B_{1}\right\rangle _{f}=\left\langle \left.  I\right\vert B_{1}\right\rangle
_{f}=f\left(  B_{1}\right) \\
f\left(  S_{f}\left(  B_{1}\right)  \right)   &  =f\left(  B_{1}^{\dagger
}\right) \\
f\left(  S_{f}^{\dagger}\left(  B_{1}\right)  \right)   &  =\left\langle
\left.  I\right\vert S_{f}^{\dagger}\left(  B_{1}\right)  \right\rangle
_{f}=\left\langle \left.  I\right\vert B_{1}\right\rangle _{f}=f\left(
B_{1}\right)  .
\end{align}


$\rfloor$

In general
\begin{equation}
\Pi_{f}\left(  \mathcal{A}\right)  \subseteq\mathcal{B}\left(  \mathcal{H}%
_{f}\right)  .
\end{equation}
Hence one can consider $S_{f}S_{f}^{\dagger}$ and $S_{f}^{\dagger}S_{f}$ which
are positive \emph{linear }operator acting on $\mathcal{A}_{f}\cong%
\mathcal{H}_{f},$ i.e $S_{f}S_{f}^{\dagger}$ and $S_{f}^{\dagger}S_{f}%
\in\mathcal{B}_{+}\left(  \mathcal{A}_{f}\right)  $

$\lceil$%
\begin{align}
S_{f}\left(  \alpha_{1}v_{1}+\alpha_{2}v_{2}\right)   &  =\overline{\alpha
}_{1}S_{f}\left(  v_{1}\right)  +\overline{\alpha}_{2}S_{f}\left(
v_{2}\right) \\
S_{f}^{\dagger}\left(  \alpha_{1}v_{1}+\alpha_{2}v_{2}\right)   &
=\overline{\alpha}_{1}S_{f}^{\dagger}\left(  v_{1}\right)  +\overline{\alpha
}_{2}S_{f}^{\dagger}\left(  v_{2}\right) \\
S_{f}S_{f}^{\dagger}\left(  \alpha_{1}v_{1}+\alpha_{2}v_{2}\right)   &
=S_{f}\left(  \overline{\alpha}_{1}S_{f}^{\dagger}\left(  v_{1}\right)
+\overline{\alpha}_{2}S_{f}^{\dagger}\left(  v_{2}\right)  \right)
=\alpha_{1}S_{f}S_{f}^{\dagger}\left(  v_{1}\right)  +\alpha_{2}S_{f}%
S_{f}^{\dagger}\left(  v_{2}\right)  .
\end{align}
$\rfloor$

that have the values
\begin{equation}
\left\langle \left.  B_{1}\right\vert \left[  S_{f}^{\dagger}S_{f}\right]
\left(  B_{2}\right)  \right\rangle _{f}=f\left(  B_{1}^{\dagger}\left[
S_{f}^{\dagger}S_{f}\right]  \left(  B_{2}\right)  \right)  \ =\overline
{f\left(  B_{2}B_{1}^{\dagger}\right)  }%
\end{equation}
and
\begin{equation}
\left\langle \left.  B_{1}\right\vert \left[  S_{f}S_{f}^{\dagger}\right]
\left(  B_{2}\right)  \right\rangle _{f}=f\left(  B_{1}^{\dagger}\left[
S_{f}S_{f}^{\dagger}\right]  \left(  B_{2}\right)  \right)  .\
\end{equation}


One can express $S_{f}$ (called the \emph{Tomita operator}) in polar form
\begin{equation}
S_{f}=\mathcal{J}\left(  S_{f}^{\dagger}S_{f}\right)  ^{\frac{1}{2}}%
\equiv\mathcal{J}\Delta^{\frac{1}{2}},
\end{equation}
or
\begin{equation}
S_{f}=\left(  S_{f}^{\dagger}S_{f}\right)  ^{-\frac{1}{2}}\mathcal{J}%
\equiv\Delta^{\frac{1}{2}}\mathcal{J},
\end{equation}
where $\mathcal{J}$ is anti unitary (thus a partial isometry) called the
\emph{modular conjugation} ($S_{f}$ is invertible as it is an involution) and
$\Delta$ is known as the \emph{modular operator}.

As $S_{f}^{\dagger}S_{f}$
\begin{equation}
S_{f}^{\dagger}S_{f}=\Delta^{\frac{1}{2}}\mathcal{J}^{\dagger}\mathcal{J}%
\Delta^{\frac{1}{2}}=\Delta=S_{f}^{\dagger}S_{f}%
\end{equation}
is a positive operator one can form the one parameter unity group,
$\mathcal{U}_{f}\left(  \mathcal{A}\right)  ,$ (by using the spectral
decomposition) called the $\emph{Modular}$ \emph{Group}
\begin{equation}
u\left(  t\right)  =\left(  S_{f}^{\dagger}S_{f}\right)  ^{it},t\in\mathbb{R}.
\end{equation}


\begin{theorem}
\begin{enumerate}
\item
\begin{equation}
\mathcal{J}\Pi_{f}\left(  \mathcal{A}\right)  ^{\prime\prime}\mathcal{J}%
=\Pi_{f}\left(  \mathcal{A}\right)  ^{\prime\prime\prime}%
\end{equation}


\item
\begin{equation}
\widehat{\Delta^{it}}\left(  \Pi_{f}\left(  \mathcal{A}\right)  ^{\prime
\prime}\right)  =\Delta^{it}\Pi_{f}\left(  \mathcal{A}\right)  ^{\prime\prime
}\Delta^{-it}=\Pi_{f}\left(  \mathcal{A}\right)  ^{\prime\prime}%
\end{equation}

\end{enumerate}
\end{theorem}

where $\widehat{\Delta^{it}}$ forms the modular group of $\Pi_{f}\left(
\mathcal{A}\right)  ^{\prime\prime}.$

For the positive convex cone $\mathcal{H}_{f+}\subset\mathcal{H}_{f}$ defined
by
\begin{equation}
\mathcal{H}_{f+}=\overline{\left\{  \left.  v\right\vert v=\left(
\mathcal{J}A\mathcal{J}A\right)  \Phi_{f}\left(  I_{\mathcal{A}}\right)
,A\in\mathcal{A}\right\}  }.
\end{equation}


The following hold:

\begin{enumerate}
\item
\begin{equation}
\mathcal{J}v=v\,\forall v\in\mathcal{H}_{f+}.
\end{equation}


\item
\begin{equation}
\left(  A\mathcal{J}A\right)  \mathcal{H}_{f+}\subset\mathcal{H}_{f+}\forall
A\in A\in\Pi_{f}\left(  \mathcal{A}\right)  ^{\prime\prime}.
\end{equation}


\item
\begin{equation}
\mathcal{J}A\mathcal{J}=A^{\dagger}\forall A\in\mathcal{A}=\Pi_{f}\left(
\mathcal{A}\right)  ^{\prime\prime}\cap\Pi_{f}\left(  \mathcal{A}\right)
^{\prime\prime\prime}=\mathcal{Z}\left(  \Pi_{f}\left(  \mathcal{A}\right)
^{\prime\prime}\right)  ,
\end{equation}
where $\mathcal{Z}\left(  \Pi_{f}\left(  \mathcal{A}\right)  ^{\prime\prime
}\right)  $ is the centre of $\Pi_{f}\left(  \mathcal{A}\right)
^{\prime\prime}.$
\end{enumerate}

\subsection{Standard Representation}

\begin{definition}
A quadruple $\left\{  \mathcal{H},\Pi,\mathcal{J},\mathcal{H}_{+}\right\}  $
is a standard representation of the $w^{\ast}$-algebra $\mathcal{A}$ if $\Pi$
: $\mathcal{A}$ $\rightarrow\mathcal{B}(\mathcal{H})$ is a representation of
$\mathcal{A}$, $\mathcal{J}$ an antiunitary involution on $\mathcal{H}$ and
$\mathcal{H}_{+}$ a self-dual cone in $\mathcal{H}$ satisfying the following conditions:
\end{definition}

\begin{problem}
\begin{enumerate}
\item $\mathcal{J}\Pi(\mathcal{A})\mathcal{J}=\Pi(\mathcal{A})^{\prime}$

\item $\mathcal{J}\Pi(A)\mathcal{J}=\Pi(A)^{\dagger}$ for all $A$ in the
center of $\mathcal{A}$

\item $\mathcal{J}\nu=\nu$ for all $\nu\in\mathcal{H}_{+}$

\item $\Pi(A)\mathcal{J}\Pi(A)\mathcal{H}_{+}\subset\mathcal{H}_{+}$ for all
$A\in\mathcal{A}$.
\end{enumerate}
\end{problem}

One of the key results in the theory of $w^{\ast}$-algebras is the following:

\begin{theorem}
\emph{Any} $w^{\ast}$-algebra $\mathcal{A}$ has a faithful standard
representation. Moreover, this representation is unique, up to unitary equivalence.
\end{theorem}

\begin{theorem}
Let $\left\{  \mathcal{H},\Pi,\mathcal{J},\mathcal{H}_{+}\right\}  $ be a
standard representation of $\mathcal{A}$. Any normal state $\omega$ on
$\mathcal{A}$ has a unique vector representative $\Phi_{\U{3c9} }\left(
I_{\mathcal{A}}\right)  $ $\in$ $\mathcal{H}_{+}$ such that $\U{3c9} (A)$
$=\left\langle \Phi_{\U{3c9} }\left(  I_{\mathcal{A}}\right)  |\Pi_{\U{3c9}
}\left(  A\right)  \Phi_{\U{3c9} }\left(  I_{\mathcal{A}}\right)
\right\rangle _{\U{3c9} }$. Moreover,%
\begin{equation}
\left\Vert \Phi_{\U{3c9} }\left(  I_{\mathcal{A}}\right)  -\Phi_{\nu}\left(
I_{\mathcal{A}}\right)  \right\Vert \leq\left\Vert \omega-\nu\right\Vert
\leq\left\Vert \Phi_{\U{3c9} }\left(  I_{\mathcal{A}}\right)  -\Phi_{\nu
}\left(  I_{\mathcal{A}}\right)  \right\Vert \left\Vert \Phi_{\U{3c9} }\left(
I_{\mathcal{A}}\right)  +\Phi_{\nu}\left(  I_{\mathcal{A}}\right)
\right\Vert
\end{equation}
holds for all normal states $\U{3c9} ,\U{3bd} $. Thus, there is a
\emph{homeomorphic correspondence} between \emph{normal states and unit
vectors} of $\mathcal{H}_{+}$. Finally,%
\begin{equation}
\Pi_{\omega}(\mathcal{A})\Phi_{\U{3c9} }\left(  I_{\mathcal{A}}\right)
=\mathcal{J}\Pi_{\omega}(\mathcal{A})^{\prime}\Phi_{\U{3c9} }\left(
I_{\mathcal{A}}\right)  ,
\end{equation}
and in particular $\left\{  \omega\text{ is faithful }\right\}
\Leftrightarrow\left\{  \Phi_{\U{3c9} }\left(  I_{\mathcal{A}}\right)  \text{
is separating for }\Pi_{\omega}(\mathcal{A})\right\}  \Leftrightarrow\left\{
\Phi_{\U{3c9} }\left(  I_{\mathcal{A}}\right)  \text{ is cyclic for }%
\Pi_{\omega}(\mathcal{A})\right\}  $.
\end{theorem}

Let $\mathcal{U}\left(  \mathcal{H}\right)  $ be the group of unitaries such
that
\begin{gather}
u\Pi_{\omega}(\mathcal{A})u^{\dagger}\in\Pi_{\omega}(\mathcal{A})\\
u\mathcal{H}_{+}\in\mathcal{H}_{+}\\
\forall u\in\mathcal{U}\left(  \mathcal{H}\right)  .
\end{gather}
Equipped with the strong operator topology $\mathcal{U}\left(  \mathcal{H}%
\right)  $ is a topological group and
\begin{equation}
\U{3c4} _{u}(A)=\Pi_{\omega}^{\text{-}1}(u\Pi_{\omega}(A)u^{\dagger})
\end{equation}
defines a continuous morphism%
\begin{equation}
\mathcal{U}\left(  \mathcal{H}\right)  \rightarrow\operatorname{Aut}\left\{
\mathcal{A}\right\}  .
\end{equation}


\begin{theorem}
The map $u\rightarrow\U{3c4} _{u}$ is a topological isomorphism. Moreover, for
any $u\in\mathcal{U}\left(  \mathcal{H}\right)  $ and \emph{any normal state}
$\U{3c9} $ on $\mathcal{A}$ one has

\begin{enumerate}
\item $\mathcal{J}$ $u$ $\mathcal{J}$ $=u$ .

\item $u\Pi_{\omega}(\mathcal{A})^{\prime}u^{\dagger}=\Pi_{\omega}%
(\mathcal{A})^{\prime}$.

\item $u^{\dagger}\Phi_{\U{3c9} }=\Phi_{\U{3c9} \circ\tau}u.$
\end{enumerate}
\end{theorem}

\begin{theorem}
In particular, if $\left\{  \mathcal{A},\U{3c4} \right\}  $ is a $w^{\ast}%
$-dynamical system then there exists a unique self-adjoint operator $L$ on
$\mathcal{H}$ such that
\begin{equation}
\Pi_{\omega}(\U{3c4} ^{t}(A))=\exp\left\{  itL\right\}  \Pi_{\omega}%
(A)\exp\left\{  -itL\right\}
\end{equation}
and
\begin{equation}
\exp\left\{  itL\right\}  \mathcal{H}_{+}\subset\mathcal{H}_{+}.
\end{equation}

\end{theorem}

\begin{definition}
The generator $L$ is called the \emph{standard Liouvillean} of the dynamical
system $\left\{  \mathcal{A},\U{3c4} \right\}  $.
\end{definition}

\begin{theorem}
The standard Liouvillean is uniquely defined up to unitary equivalence. If
$\U{3c9} $ is a modular $\U{3c4} $ -invariant state then the induced GNS
representation is standard and the $\U{3c9} $-Liouvillean coincide with the
standard Liouvillean. In particular this the case if $\U{3c9} $ is a
\emph{KMS} state for $\U{3c4} $ .
\end{theorem}

\begin{theorem}
Let $L$ be the standard Liouvillean of the $w^{\ast}$-dynamical system
$\left\{  \mathcal{A},\U{3c4} \right\}  $. Then

\begin{enumerate}
\item L has no eigenvalues if and only if there is no normal $\U{3c4} $
-invariant state on $\mathcal{A}$,

\item $\ker(L)$ is one-dimensional if and only if there is a unique normal
$\U{3c4} $ -invariant state $\U{3c9} $ on $\mathcal{A}$. In this case
$\Phi_{\U{3c9} }\left(  I_{\mathcal{A}}\right)  $ is the unique unit vector in
$\ker(L)\cap$ $\mathcal{H}_{+}$.
\end{enumerate}
\end{theorem}

\section{The finite dimensional case}

It is instructive to work out the standard representation of a finite
dimensional von Neumann algebra $\Pi_{\U{3c9} }\left(  \mathcal{A}\right)  $
$\subset$ $\mathcal{B}(\mathbb{C}^{n})$, $n<\infty,$ for any state $\omega$.
This case is particularly simple since $\mathcal{B}$($\mathbb{C}^{n}$ ) is
itself a Hilbert space (Hilbert-Schmidt) for the inner product $\left\langle
X|Y\right\rangle $ $=$ $\operatorname{Tr}\left\{  X^{\dagger}Y\right\}  $. One
has
\begin{equation}
\mathcal{B}(\mathbb{C}^{n})=\Phi_{\U{3c9} }\left(  \mathcal{A}\right)
\oplus\Phi_{\U{3c9} }\left(  \mathcal{A}\right)  ^{\perp}%
\end{equation}
and the predual $\Pi_{\U{3c9} }\left(  \mathcal{A}\right)  _{\ast}$ can be
identified with $\Pi_{\U{3c9} }\left(  \mathcal{A}\right)  $.

One can define a mixed faithful normal state
\begin{equation}
\U{3c9} _{\U{3c1} }:\mathcal{A\rightarrow}\mathbb{R}%
\end{equation}
by
\begin{equation}
\U{3c9} _{\U{3c1} }\left(  A\right)  =\left\langle \Phi_{\U{3c9} _{\U{3c1} }%
}\left(  I_{\mathcal{A}}\right)  |\Pi(A)\Phi_{\U{3c9} _{\U{3c1} }}\left(
I_{\mathcal{A}}\right)  \right\rangle =\operatorname{Tr}\left\{  A\U{3c1}
\right\}  ,A\in\mathcal{A}%
\end{equation}
where
\begin{equation}
\U{3c1} =\frac{1}{n}\sum_{1\leq j\leq n}\U{3c1} _{j}%
\end{equation}
and the component states have the properties
\begin{align}
\U{3c1} _{j}  &  =\left\vert e_{j}\right\rangle \left\langle e_{j}\right\vert
\\
\left[  \U{3c1} _{j},\U{3c1} _{k}\right]   &  =0\\
\left\{  \U{3c1} _{j},\U{3c1} _{k}\right\}   &  =2\delta_{jk}\\
\U{3c1} _{j}^{2}  &  =\U{3c1} _{j}\\
\U{3c1} _{j}  &  \geq0\\
\operatorname{Tr}\left\{  \U{3c1} _{j}\right\}   &  =1\\
\mathbb{C}^{n}  &  =\operatorname*{linspan}\left\{  e_{j},1\leq j\leq
n\right\}  .
\end{align}
$\lceil$ One can generalize to
\begin{align}
\U{3c1}  &  =\frac{1}{n}\sum_{1\leq j\leq n}\beta_{j}\left\vert e_{j}%
\right\rangle \left\langle e_{j}\right\vert \\
\sum_{1\leq j\leq n}\beta_{j}  &  =n,\beta_{j}>0,1\leq j\leq n
\end{align}
to produce a more general faithful state $\U{3c1} \ $which is a general
\emph{definite }convex combination of $\left\{  \left\vert e_{j}\right\rangle
\left\langle e_{j}\right\vert \right\}  $
\begin{align}
\U{3c1}  &  =\sum_{1\leq j\leq n}\frac{\beta_{j}}{n}\left\vert e_{j}%
\right\rangle \left\langle e_{j}\right\vert =\sum_{1\leq j\leq n}\alpha
_{j}\left\vert e_{j}\right\rangle \left\langle e_{j}\right\vert \\
\sum_{1\leq j\leq n}\alpha_{j}  &  =1,1\geq\alpha_{j}>0,1\leq j\leq n.
\end{align}
It should be noted that such $\U{3c1} $ are invertible.$\rfloor$

The state $\U{3c1} $ is obviously faithful as
\begin{equation}
\sum_{1\leq j\leq n}\U{3c1} _{j}=I_{n}.
\end{equation}
As $\U{3c1} \geq0$
\begin{equation}
\U{3c1} =QQ^{\dagger}%
\end{equation}
and one can consider the positive square root
\begin{equation}
\left\vert Q\right\vert =\left(  QQ^{\dagger}\right)  ^{\frac{1}{2}}=\U{3c1}
^{\frac{1}{2}}%
\end{equation}
so that
\begin{equation}
\U{3c1} =\left\vert Q\right\vert ^{2}.
\end{equation}
The square root of $\U{3c1} $ is given by
\begin{equation}
\U{3c1} ^{\frac{1}{2}}=\frac{1}{\sqrt{n}}\sum_{1\leq j\leq n}\U{3c1} _{j}.
\end{equation}
One should note that
\begin{equation}
\operatorname{Tr}\left\{  \U{3c1} ^{\frac{1}{2}}\right\}  \neq1
\end{equation}
and thus $\U{3c1} ^{\frac{1}{2}}$ is not a state, \emph{though it is a
normalized vector in the Hilbert Schmidt space} $\mathcal{B}_{2}%
(\mathbb{C}^{n}).$

The vector space dimension of $\mathcal{B}_{2}(\mathbb{C}^{n})$ is $n^{2},$
while the space of bounded linear operators, $\mathcal{B}\left(
\mathcal{B}_{2}(\mathbb{C}^{n})\right)  ,$ acting in $\mathcal{B}%
_{2}(\mathbb{C}^{n})$ has vector space dimension $n^{4}.$ Expectation values
in the vector state $\U{3c1} ^{\frac{1}{2}}$ $\in\mathcal{B}_{2}%
(\mathbb{C}^{n})$ of operators, $\mathfrak{A},$ belonging to $\mathcal{B}%
\left(  \mathcal{B}_{2}(\mathbb{C}^{n})\right)  $ are given by
\begin{equation}
\left\langle \left.  \U{3c1} ^{\frac{1}{2}}\right\vert \mathfrak{A}\U{3c1}
^{\frac{1}{2}}\right\rangle _{HS}=\operatorname{Tr}\left\{  \U{3c1} ^{\frac
{1}{2}}\mathfrak{A}\left(  \U{3c1} ^{\frac{1}{2}}\right)  \right\}  .
\end{equation}
If one restricts attention to the subspace of left multiplication operators
generated $\mathcal{B}(\mathbb{C}^{n})$ i.e.
\begin{equation}
A_{L}\left(  B\right)  =AB\,\forall B\in\mathcal{B}(\mathbb{C}^{n}),
\end{equation}
the expectaion values are given by
\begin{equation}
\left\langle \left.  \U{3c1} ^{\frac{1}{2}}\right\vert A_{L}\left(  \U{3c1}
^{\frac{1}{2}}\right)  \right\rangle _{HS}=\operatorname{Tr}\left\{  \U{3c1}
^{\frac{1}{2}}A\U{3c1} ^{\frac{1}{2}}\right\}  =\operatorname{Tr}\left\{
A\U{3c1} \right\}  ,
\end{equation}
which are equal to the expectation values of $\mathcal{B}(\mathbb{C}^{n})$ in
the state $\U{3c9} _{\U{3c1} }$ i.e.
\begin{equation}
\operatorname{Tr}\left\{  A\U{3c1} \right\}  =\U{3c9} _{\U{3c1} }\left(
A\right)  .
\end{equation}


The GNS Hilbert space, $\mathcal{H}_{\U{3c9} _{\U{3c1} }},$ generated by the
state $\U{3c9} _{\U{3c1} }$ is given by
\begin{equation}
\mathcal{H}_{\U{3c9} _{\U{3c1} }}=\Phi_{\U{3c9} _{\U{3c1} }}\left(
\mathcal{B}(\mathbb{C}^{n})\right)  =\left\{  \left.  \Phi_{\U{3c9} _{\U{3c1}
}}\left(  A\right)  \right\vert \operatorname{Tr}\left\{  A^{\dagger}A\U{3c1}
\right\}  \neq0\right\}  =\mathcal{B}(\mathbb{C}^{n}),
\end{equation}
when equipped with the inner product
\begin{equation}
\left\langle \Phi_{\U{3c9} _{\U{3c1} }}\left(  A\right)  |\Phi_{\U{3c9}
_{\U{3c1} }}\left(  B\right)  \right\rangle _{\U{3c9} _{\U{3c1} }%
}=\operatorname{Tr}\left\{  A^{\dagger}B\U{3c1} \right\}
\end{equation}
as the state $\U{3c9} _{\U{3c1} }$ is faithful. The relationship between the
HS\ norm
\begin{equation}
\left\Vert A\right\Vert _{HS}=\left(  \operatorname{Tr}\left\{  A^{\dagger
}A\right\}  \right)  ^{\frac{1}{2}}%
\end{equation}
and norm of $\mathcal{H}_{\U{3c9} _{\U{3c1} }}$
\begin{equation}
\left\Vert A\right\Vert _{\mathcal{H}_{\U{3c9} _{\U{3c1} }}}=\left(
\operatorname{Tr}\left\{  A^{\dagger}A\U{3c1} \right\}  \right)  ^{\frac{1}%
{2}}%
\end{equation}
is
\begin{equation}
\left\Vert A\right\Vert _{\mathcal{H}_{\U{3c9} _{\U{3c1} }}}\leq\left\Vert
A\right\Vert _{HS},
\end{equation}
thus in infinite dimensions the closure of the Hilbert-Schmidt space is smaller.

The GNS\ representation, $\Pi_{\U{3c9} _{\U{3c1} }},$ generated by the state
$\U{3c9} _{\U{3c1} }$ is given by
\begin{equation}
\Pi_{\U{3c9} _{\U{3c1} }}\left(  A\right)  \Phi_{\U{3c9} _{\U{3c1} }}\left(
B\right)  =\Phi_{\U{3c9} _{\U{3c1} }}\left(  AB\right)  \,\forall
B\in\mathcal{B}(\mathbb{C}^{n}).
\end{equation}
As $\Pi_{\U{3c9} _{\U{3c1} }}$ is faithful and thus a $c^{\ast}$-algebra
isomorphism $\Pi_{\U{3c9} _{\U{3c1} }}^{\text{-}1}$ exists.

The expectation values wrt the state $\U{3c9} _{\U{3c1} }$ can be expressed in
terms of the GNS\ representation or the Hilbert Schmidt space as
\begin{equation}
\operatorname{Tr}\left\{  A\U{3c1} \right\}  =\left\langle \Phi_{\U{3c9}
_{\U{3c1} }}\left(  I_{n}\right)  |\Pi_{\U{3c9} _{\U{3c1} }}\left(  A\right)
\Phi_{\U{3c9} _{\U{3c1} }}\left(  I_{n}\right)  \right\rangle _{\U{3c9}
_{\U{3c1} }}=\left\langle \left.  \U{3c1} ^{\frac{1}{2}}\right\vert A\U{3c1}
^{\frac{1}{2}}\right\rangle _{HS}.
\end{equation}
A basis of $\mathcal{B}(\mathbb{C}^{n})$ is given by $\left\{  \left.
\left\vert e_{j}\right\rangle \left\langle e_{k}\right\vert \right\vert 1\leq
j,k\leq n\right\}  ,$where $\left\{  \left.  \left\vert e_{j}\right\rangle
\right\vert 1\leq j\leq n\right\}  $ is a conb for $\mathbb{C}^{n}.$ As
$\mathbb{C}^{n}$ is finite dimensional
\begin{equation}
\mathcal{B}(\mathbb{C}^{n})=\mathcal{B}_{0}(\mathbb{C}^{n})=\mathcal{B}%
_{1}(\mathbb{C}^{n})=\mathcal{B}_{2}(\mathbb{C}^{n}).
\end{equation}
Thus the basis $\left\{  \left.  \left\vert e_{j}\right\rangle \left\langle
e_{k}\right\vert \right\vert 1\leq j,k\leq n\right\}  $ is a conb for
$\mathcal{B}(\mathbb{C}^{n})$ wrt the HS inner product
\begin{equation}
\left\langle \left.  A\right\vert B\right\rangle _{HS}=\operatorname{Tr}%
\left\{  A^{\dagger}B\right\}
\end{equation}
i.e.
\begin{equation}
\left\langle \left.  \left\vert e_{j}\right\rangle \left\langle e_{k}%
\right\vert \right\vert \left\vert e_{l}\right\rangle \left\langle
e_{m}\right\vert \right\rangle _{HS}=\operatorname{Tr}\left\{  \left\vert
e_{k}\right\rangle \left\langle e_{j}\right\vert \left\vert e_{l}\right\rangle
\left\langle e_{m}\right\vert \right\}  =\delta_{jl}\delta_{km}=\delta
_{\left(  jk\right)  \left(  lm\right)  }.
\end{equation}
The basis $\left\{  \left.  \Phi_{\U{3c9} _{\U{3c1} }}\left(  \frac{1}%
{\sqrt{n}}\left\vert e_{j}\right\rangle \left\langle e_{k}\right\vert \right)
\right\vert 1\leq j,k\leq n\right\}  $ is conb for $\Phi_{\U{3c9} _{\U{3c1} }%
}\left(  \mathcal{B}(\mathbb{C}^{n})\right)  $ wrt the weighted HS inner
product
\begin{align}
&  \left\langle \Phi_{\U{3c9} _{\U{3c1} }}\left(  \sqrt{n}\left\vert
e_{j}\right\rangle \left\langle e_{k}\right\vert \right)  |\Phi_{\U{3c9}
_{\U{3c1} }}\sqrt{n}\left(  \left\vert e_{l}\right\rangle \left\langle
e_{m}\right\vert \right)  \right\rangle _{\U{3c9} _{\U{3c1} }}%
=n\operatorname{Tr}\left\{  \left\vert e_{k}\right\rangle \left\langle
e_{j}\right\vert \left\vert e_{l}\right\rangle \left\langle e_{m}\right\vert
\U{3c1} \right\}  =\delta_{jl}n\operatorname{Tr}\left\{  \left\vert
e_{k}\right\rangle \left\langle e_{m}\right\vert \U{3c1} \right\}
=\delta_{jl}n\left\langle \left.  e_{m}\right\vert \U{3c1} e_{k}\right\rangle
_{\mathbb{C}^{n}}\\
&  =\delta_{jl}n\left\langle e_{m}\left\vert \frac{1}{n}\sum_{1\leq p\leq
n}\left\vert e_{p}\right\rangle \left\langle e_{p}\right\vert \right.
e_{k}\right\rangle _{\mathbb{C}^{n}}=\delta_{jl}\delta_{km}=\delta_{\left(
jk\right)  \left(  lm\right)  }.
\end{align}


The adjoint map in $\mathcal{B}(\mathbb{C}^{n})$ is given by
\begin{equation}
\left\langle \left.  A^{\dagger}v_{1}\right\vert v_{2}\right\rangle
_{\mathbb{C}^{n}}=\left\langle \left.  v_{1}\right\vert Av_{2}\right\rangle
_{\mathbb{C}^{n}};\,\forall v_{2}\in\mathbb{C}^{n}%
\end{equation}
The $\dagger$ operation in $\mathcal{B}(\mathbb{C}^{n})$ is mapped to
$\dagger$ operation in $\Pi_{\U{3c9} _{\U{3c1} }}\left(  \mathcal{B}%
(\mathbb{C}^{n})\right)  $ as for all $X,Y\in\mathcal{B}(\mathbb{C}^{n})$
\begin{align}
\left\langle \Phi_{\U{3c9} _{\U{3c1} }}\left(  X\right)  |\Pi_{\U{3c9}
_{\U{3c1} }}\left(  A\right)  \Phi_{\U{3c9} _{\U{3c1} }}\left(  Y\right)
\right\rangle _{\U{3c9} _{\U{3c1} }}  &  =\operatorname{Tr}\left\{
X^{\dagger}AY\U{3c1} \right\} \\
\left\langle \Phi_{\U{3c9} _{\U{3c1} }}\left(  X\right)  |\Pi_{\U{3c9}
_{\U{3c1} }}\left(  A^{\dagger}\right)  \Phi_{\U{3c9} _{\U{3c1} }}\left(
Y\right)  \right\rangle _{\U{3c9} _{\U{3c1} }}  &  =\operatorname{Tr}\left\{
X^{\dagger}A^{\dagger}Y\U{3c1} \right\} \\
\left\langle \Phi_{\U{3c9} _{\U{3c1} }}\left(  X\right)  |\Pi_{\U{3c9}
_{\U{3c1} }}\left(  A\right)  ^{\dagger}\Phi_{\U{3c9} _{\U{3c1} }}\left(
Y\right)  \right\rangle _{\U{3c9} _{\U{3c1} }}  &  =\left\langle \Pi_{\U{3c9}
_{\U{3c1} }}\left(  A\right)  \Phi_{\U{3c9} _{\U{3c1} }}\left(  X\right)
|\Phi_{\U{3c9} _{\U{3c1} }}\left(  Y\right)  \right\rangle _{\U{3c9} _{\U{3c1}
}}=\left\langle \Phi_{\U{3c9} _{\U{3c1} }}\left(  AX\right)  |\Phi_{\U{3c9}
_{\U{3c1} }}\left(  Y\right)  \right\rangle _{\U{3c9} _{\U{3c1} }}\\
&  =\operatorname{Tr}\left\{  \left(  AX\right)  ^{\dagger}Y\U{3c1} \right\}
=\operatorname{Tr}\left\{  X^{\dagger}A^{\dagger}Y\U{3c1} \right\}
=\left\langle \Phi_{\U{3c9} _{\U{3c1} }}\left(  X\right)  |\Pi_{\U{3c9}
_{\U{3c1} }}\left(  A^{\dagger}\right)  \Phi_{\U{3c9} _{\U{3c1} }}\left(
Y\right)  \right\rangle _{\U{3c9} _{\U{3c1} }}%
\end{align}
thus
\begin{equation}
\Pi_{\U{3c9} _{\U{3c1} }}\left(  A\right)  ^{\dagger}=\Pi_{\U{3c9} _{\U{3c1}
}}\left(  A^{\dagger}\right)
\end{equation}
for any state $\U{3c9} _{\U{3c1} }.$ This also follows from $\Pi_{\U{3c9}
_{\U{3c1} }}$ being a $c^{\ast}$-algebra morphism.

Consider $\mathcal{H}_{\U{3c9} _{\U{3c1} }}$ $=\Phi_{\U{3c9} _{\U{3c1} }%
}\left(  \mathcal{A}\right)  $ as a Hilbert space $\subseteq$ $\mathcal{B}%
_{2}(\mathbb{C}^{n})$. Then $\Phi_{\U{3c9} _{\U{3c1} }}\left(  I_{\mathcal{A}%
}\right)  $ is a unit vector in $\mathcal{H}_{\U{3c9} _{\U{3c1} }}$. Moreover,
the map
\begin{equation}
\Pi_{\U{3c9} _{\U{3c1} }}:\mathcal{A}\rightarrow\Pi_{\U{3c9} _{\U{3c1} }%
}\left(  \mathcal{A}\right)  \subseteq\mathcal{B(H)}%
\end{equation}
defined by
\begin{equation}
\Pi_{\U{3c9} _{\U{3c1} }}(A)\Phi_{\U{3c9} _{\U{3c1} }}\left(  X\right)
=\Phi_{\U{3c9} _{\U{3c1} }}\left(  AX\right)
\end{equation}
is a $\ast$-morphism such that
\begin{equation}
\U{3c9} _{\U{3c1} }(A)=\left\langle \Phi_{\U{3c9} _{\U{3c1} }}\left(
I_{\mathcal{A}}\right)  |\Pi_{\U{3c9} _{\U{3c1} }}(A)\Phi_{\U{3c9} _{\U{3c1}
}}\left(  I_{\mathcal{A}}\right)  \right\rangle _{\U{3c9} _{\U{3c1} }%
}=\left\langle \left.  \U{3c1} ^{\frac{1}{2}}\right\vert A\U{3c1} ^{\frac
{1}{2}}\right\rangle _{HS}.
\end{equation}
Denote by $P_{\U{3c9} _{\U{3c1} }}$ the orthogonal projection on $\ker\left\{
\U{3c9} _{\U{3c1} }\right\}  ^{\perp}$. Clearly $P_{\U{3c9} _{\U{3c1} }}%
\in\mathcal{A}$ and since $\U{3c1} _{j}\geq0$ one has%
\begin{equation}
\ker\left\{  \Phi_{\U{3c9} _{\U{3c1} }}\right\}  \subset%
%TCIMACRO{\tbigcap \limits_{1\leq j\leq n}}%
%BeginExpansion
{\textstyle\bigcap\limits_{1\leq j\leq n}}
%EndExpansion
\ker\left\{  \Phi_{\U{3c1} _{j}}\right\}
\end{equation}
and hence
\begin{equation}
\ker\left\{  \Phi_{\U{3c9} _{\U{3c1} }}\right\}  \subset\ker\left\{
A\right\}  \forall A\in\mathcal{A}.
\end{equation}
It follows that
\begin{equation}
\operatorname{Ran}\left\{  A\right\}  \subset\operatorname{Ran}\left\{
P_{\U{3c9} _{\U{3c1} }}\right\}
\end{equation}
and hence
\begin{equation}
A=P_{\U{3c9} _{\U{3c1} }}A\forall A\in\mathcal{A}.
\end{equation}
Since the last identity is equivalent to
\begin{equation}
A^{\dagger}=A^{\dagger}P_{\U{3c9} _{\U{3c1} }}%
\end{equation}
we conclude that
\begin{equation}
A=AP_{\U{3c9} _{\U{3c1} }}=P_{\U{3c9} _{\U{3c1} }}A=P_{\U{3c9} _{\U{3c1} }%
}AP_{\U{3c9} _{\U{3c1} }}\forall A\in\mathcal{A}%
\end{equation}
i.e., that $P_{\U{3c9} _{\U{3c1} }}$ \emph{is the unit of} $\mathcal{A}$.
Since there exists $T\in\mathcal{A}$ such that $T\rho^{\frac{1}{2}}=P_{\U{3c9}
_{\U{3c1} }}$ we can write
\begin{equation}
\Pi_{\U{3c9} _{\U{3c1} }}(XT)\rho^{\frac{1}{2}}=X
\end{equation}
and conclude that
\begin{equation}
\Pi_{\U{3c9} _{\U{3c1} }}(\mathcal{A})\Phi_{\U{3c9} _{\U{3c1} }}\left(
I_{\mathcal{A}}\right)  =\mathcal{H}_{\U{3c9} _{\U{3c1} }}.
\end{equation}
We have shown that $\left\{  \mathcal{H}_{\U{3c9} _{\U{3c1} }},\Pi_{\U{3c9}
_{\U{3c1} }},\Phi_{\U{3c9} _{\U{3c1} }}\left(  I_{\mathcal{A}}\right)
\right\}  $ is the GNS representation of $\mathcal{A}$ induced by $\U{3c9}
_{\U{3c1} }$. Note that since $\U{3c9} _{\U{3c1} }$ is a normal faithful
state, this representation is itself faithful.

The formulas
\begin{align}
\mathcal{J}\left(  X\right)   &  =X^{\dagger}\\
\U{2206} _{\U{3c9} _{\U{3c1} }}^{\frac{1}{2}}XT  &  =\rho^{\frac{1}{2}}XT
\end{align}
define an anti-unitary involution and a positive self-adjoint operator on
$\mathcal{H}_{\U{3c9} _{\U{3c1} }}$ such that%
\begin{equation}
\mathcal{J}\U{2206} _{\U{3c9} _{\U{3c1} }}^{\frac{1}{2}}\Pi_{\U{3c9} _{\U{3c1}
}}\left(  A\right)  \Phi_{\U{3c9} _{\U{3c1} }}\left(  I_{\mathcal{A}}\right)
=\Pi_{\U{3c9} _{\U{3c1} }}\left(  A\right)  ^{\dagger}\Phi_{\U{3c9} _{\U{3c1}
}}\left(  I_{\mathcal{A}}\right)  .
\end{equation}
Thus, $\mathcal{J}$ and $\U{2206} _{\U{3c9} _{\U{3c1} }}^{\frac{1}{2}}$ are
the modular conjugation and the modular operator of the pair $\left\{
\Pi_{\U{3c9} _{\U{3c1} }}(\mathcal{A}),\Phi_{\U{3c9} _{\U{3c1} }}\left(
I_{\mathcal{A}}\right)  \right\}  $.

Elements of the natural cone are given by
\begin{equation}
\Pi_{\U{3c9} _{\U{3c1} }}\left(  A\right)  \mathcal{J}\Pi_{\U{3c9} _{\U{3c1}
}}\left(  A\right)  \Phi_{\U{3c9} _{\U{3c1} }}\left(  I_{\mathcal{A}}\right)
=\Pi_{\U{3c9} _{\U{3c1} }}\left(  A\right)  \rho^{\frac{1}{2}}\Pi_{\U{3c9}
_{\U{3c1} }}\left(  A\right)  ^{\dagger}%
\end{equation}
from which we can conclude that%
\begin{equation}
\mathcal{H}_{+}=\left\{  A\in\mathcal{A}\text{%
%TCIMACRO{\TEXTsymbol{\vert}}%
%BeginExpansion
$\vert$%
%EndExpansion
}A\geq0\right\}  .
\end{equation}


Let $c$ be an element of $\Pi_{\U{3c9} _{\U{3c1} }}(\mathcal{A})^{\prime}$.
For all $A\in\Pi_{\U{3c9} _{\U{3c1} }}\left(  A\right)  $ and $X\in
\mathcal{H}_{\U{3c9} _{\U{3c1} }}$ one has
\begin{equation}
cA\left(  X\right)  =c\Pi_{\U{3c9} _{\U{3c1} }}(A)\left(  X\right)
=\Pi_{\U{3c9} _{\U{3c1} }}(\mathcal{A})c(X)=Ac\left(  X\right)  .
\end{equation}


Setting $X=P_{\U{3c9} _{\U{3c1} }}$ , the unit of $\mathcal{A}$, and
\begin{equation}
B=cP_{\U{3c9} _{\U{3c1} }}\in\mathcal{A},
\end{equation}
we get $cA=AB$. We conclude that $\Pi_{\U{3c9} _{\U{3c1} }}(\mathcal{A}%
)^{\prime}$ consists of the linear maps%
\begin{equation}
X\rightarrow XB
\end{equation}
with $B\in\mathcal{A}$. Thus,
\begin{equation}
\mathcal{J\Pi_{\U{3c9} _{\U{3c1} }}(\mathcal{A})J}=\Pi_{\U{3c9} _{\U{3c1} }%
}(\mathcal{A})^{\prime}%
\end{equation}
and we have obtained the standard representation of the finite dimensional von
Neumann algebra $\mathcal{A}$.

Any normal state $\U{3bd} $ on $\mathcal{A}$ is given by $\U{3bd}
(A)=\operatorname{Tr}\left\{  A\U{3c1} \right\}  $ for a density matrix
$\U{3c1} \in\mathcal{A}$. It follows that%
\begin{equation}
\U{3bd} \rightarrow\Phi_{\U{3c9} _{\U{3c1} }}\left(  \U{3bd} \right)  =\U{3c1}
^{\frac{1}{2}}\in\mathcal{H}_{+}%
\end{equation}
is the homeomorphism described in the preceding Theorem.

\section{Decomposition}

Any $A$ $\in\mathcal{A}$ can be written as a linear combination of 4
non-negative elements of $\mathcal{A}$:

First we can decompose into the sum of two self adjoint operators
\begin{align}
A &  =A_{R}+iA_{I}\\
A_{R} &  =\frac{1}{2}\left(  A+A^{\dagger}\right)  \\
A_{I} &  =-\frac{i}{2}\left(  A-A^{\dagger}\right)  ,
\end{align}
which then can be decomposed into the difference of two positive operators%
\begin{align}
A_{R} &  =A_{Rp}-A_{m}\\
A_{Rp} &  =\frac{1}{2}\left(  \left\vert A_{R}\right\vert +A_{R}\right)
=\frac{1}{4}\left(  \left\vert \left(  A+A^{\dagger}\right)  \right\vert
+\left(  A+A^{\dagger}\right)  \right)  \\
A_{R_{m}--} &  =\frac{1}{2}\left(  \left\vert A_{R}\right\vert -A_{R}\right)
=\frac{1}{4}\left(  \left\vert \left(  A+A^{\dagger}\right)  \right\vert
-\left(  A+A^{\dagger}\right)  \right)
\end{align}%
\begin{align}
A_{I} &  =A_{I+}-A_{I\text{\textit{-}}-}\\
A_{I+} &  =\frac{1}{2}\left(  \left\vert A_{I}\right\vert +A_{I}\right)
=\frac{1}{4}\left(  \left\vert \left(  A-A^{\dagger}\right)  \right\vert
-i\left(  A-A^{\dagger}\right)  \right)  \\
A_{I\text{\textit{-}}-} &  =\frac{1}{2}\left(  \left\vert A_{I}\right\vert
-A_{I}\right)  =\frac{1}{4}\left(  \left\vert \left(  A-A^{\dagger}\right)
\right\vert +i\left(  A-A^{\dagger}\right)  \right)  ,
\end{align}
where one can also show that
\begin{align}
A_{R+}A_{R\text{-}--} &  =0\\
A_{I+}A_{I\text{\textit{-}}-} &  =0.
\end{align}%
\begin{equation}
A=A_{R+}-A_{R\text{-}}+i\left(  A_{I+}-A_{I\text{\textit{-}}}\right)
=\frac{1}{2}\left\{  \left(  \left\vert A_{R}\right\vert +A_{R}\right)
-\left(  \left\vert A_{R}\right\vert -A_{R}\right)  +i\left(  \left(
\left\vert A_{I}\right\vert +A_{I}\right)  -\left(  \left\vert A_{I}%
\right\vert -A_{I}\right)  \right)  \right\}  .
\end{equation}



\end{document}